\section{Marked cubic surfaces and outer Galois points}

The trace-code theorem becomes substantially more concrete in the minimal
norm-six shell.

\subsection{Unique decomposition of a norm-six trace}

Every $v\in E_8$ with $v^2=6$ has a unique unordered decomposition
\[
\boxed{v=\alpha+\beta}
\]
where $\alpha$ and $\beta$ are roots with
\[
\alpha^2=\beta^2=2,
\qquad
\langle\alpha,\beta\rangle=1.
\]
Indeed, a root $\alpha$ occurs precisely when $\langle v,\alpha\rangle=3$,
for then $v-\alpha$ is again a root.  There can be at most two such roots,
and they are complementary.  Counting gives the same result: each of the
$240$ roots has $56$ roots of inner product one, so the number of unordered
pairs is
\[
\frac{240\cdot56}{2}=6720,
\]
exactly the size of the norm-six shell.

Let $S_\alpha,S_\beta$ be the corresponding root sections.  They are pairwise
disjoint from each other and from $O$, because
\[
S_\alpha\cdot S_\beta
=1-\langle\alpha,\beta\rangle=0.
\]
The minimal trisection divisor satisfies
\[
\boxed{
D_v=F+O+S_\alpha+S_\beta.
}
\]

\subsection{Anticanonical cubic-surface model}

Contract the three disjoint $(-1)$-curves
\[
O,\quad S_\alpha,\quad S_\beta
\]
to obtain $\varphi:S\to X$.  Since $K_S=-F$,
\[
K_S=\varphi^*K_X+O+S_\alpha+S_\beta
\]
and hence
\[
\boxed{D_v=\varphi^*(-K_X).}
\]
The surface $X$ is a smooth degree-three del Pezzo surface.  Its
anticanonical system embeds it as a smooth cubic surface in $\mathbf P^3$, and
$|D_v|$ is the pullback of the hyperplane system.

Let $p_0,p_1,p_2$ be the images of the three contracted sections.  The
elliptic pencil $|F|$ is the strict transform of anticanonical hyperplanes
through all three points.  Since this system is one-dimensional, the points
are collinear.  Their line $\ell$ is not contained in $X$, and
\[
\ell\cap X=p_0+p_1+p_2.
\]

For a hyperplane $H$ not containing $\ell$, the curve
\[
C_H=X\cap H
\]
is a plane cubic and the induced trisection map is projection from
\[
q_H=H\cap\ell.
\]
Thus the Riemann--Roch search is equivalent to a marked cubic-surface problem.

\subsection{Cyclicity is an outer Galois-point condition}

The degree-three projection $C_H\to\mathbf P^1$ is cyclic exactly when $q_H$
is an outer Galois point of the plane cubic.  Over an algebraic closure, one
can then choose coordinates with
\[
q_H=(1:0:0),
\qquad
C_H:\quad X^3+G_3(Y,Z)=0.
\]
The deck transformation is $X\mapsto\zeta_3X$.  In particular,
\[
\boxed{j(C_H)=0.}
\]
Conversely, the displayed outer-Galois form gives a cyclic cubic trisection.
Descent to $\Q$ and positive rational rank remain separate arithmetic tests.

The new exact search is therefore:
\begin{enumerate}[leftmargin=*]
\item recover the unique root pair for each trace class;
\item construct the marked cubic surface $(X,\ell)$;
\item solve for hyperplanes whose moving projection centre is an outer Galois point;
\item retain smooth $j=0$ curves with positive rational rank;
\item group distinct trace classes producing the same cover and compute their Hermitian rank.
\end{enumerate}

The first decisive milestone is one common cyclic cover arising from two
distinct trace classes.  The final target is eleven Hermitian-independent
classes over one positive-rank cover.
